\documentclass[11pt,a4paper]{article}
\usepackage[margin=2cm]{geometry}
\usepackage{amsmath}
\usepackage{amssymb}
\usepackage{graphicx}
\usepackage{algorithm}
\usepackage{algpseudocode}
\usepackage{booktabs}
\usepackage{float}
\usepackage{indentfirst}

\usepackage{titlesec}
\titlespacing*{\section}{0pt}{4pt}{2pt}

\algrenewcommand\algorithmicrequire{\textbf{Input:}}
\algrenewcommand\algorithmicensure{\textbf{Output:}}

\title{\vspace{-0.5cm}\textbf{Batch Bayesian Optimization for Bioprocess Parameter Tuning}}
\author{Hammerschmidt, Jakob Elias \and Al Hachem, Marc \and Elmehelmy, Seif Ahmed Moheb}
\date{}

\begin{document}
\maketitle
\vspace{-0.5cm}

\section{Big Picture and Intuition}

Our algorithm optimizes bioreactor conditions (temperature, pH, three feed rates, and cell type) to maximize product titre under strict constraints: 6 initial samples, 15 batches of 5 experiments, and a 60-second runtime.
The core principle is Bayesian optimization, which builds a probabilistic surrogate model to guide sampling toward promising regions while quantifying uncertainty to balance exploration and exploitation.

We use a Gaussian Process (GP) as the surrogate, predicting both expected titre and confidence intervals at any configuration.
The Upper Confidence Bound (UCB) acquisition function selects points with high predicted mean or high uncertainty, naturally trading off exploitation of known good regions with exploration of uncertain ones.
For batch selection, we add a repulsion penalty preventing clustered selections, plus a cell-type diversity bonus ensuring multiple cell types are tested per batch.

A key challenge is avoiding local optima, since the objective landscape may have multiple peaks and early convergence to a suboptimal region wastes experimental budget.
We address this through three mechanisms: adaptive hyperparameters that start with broad exploration and progressively tighten; local search with cross-cell-type sampling that tests whether nearby cell types outperform the current best; and automatic exploration boosts when improvement stalls or one cell type dominates recent samples.
The mixed kernel GP handles both continuous variables (via RBF kernel) and the categorical cell type (via multiplicative kernel with $\theta_{\text{cat}}=0.25$), enabling knowledge transfer between cell types while respecting their distinct responses.

\section{Methodology}

\textbf{Initial Design and Candidate Generation.}
We initialize with 6 points from a 5D Sobol sequence, cycling cell types to ensure 2 samples per type.
Sobol sequences provide quasi-random, low-discrepancy coverage superior to random sampling.
For candidate generation, we use 6D Sobol (5 continuous + 1 categorical) with counts of 12,000 for iterations 1--3, 8,000 for iterations 4--6, and 5,000 thereafter.
Candidates are shuffled to prevent structured bias.

\textbf{Local Search with Cell-Type Variation.}
From iteration 3 onward, we augment global candidates with 2,500--4,000 local samples in a 12\% window around the current best, with a minimum window size of 3\% per dimension to avoid degenerate local boxes when ranges are small.
Only 70\% use the best cell type; the remaining 30\% test other cell types in the same continuous region, helping discover whether a different cell type performs better at similar conditions.

\textbf{Gaussian Process Model.}
The GP uses a product kernel $K = K_{\text{RBF}} \times K_{\text{cat}}$, where the RBF component handles continuous variables and the categorical component returns 1.0 for matching cell types, $\theta_{\text{cat}}=0.25$ otherwise.
Length scales are adaptive: starting at 35\% of variable range (broad) and shrinking to 20\% as data accumulates (narrow), promoting early exploration then late refinement.
Outputs are normalized and Cholesky decomposition enables $O(n^2)$ predictions.
Noise is adaptive and estimated using a normalized range-based heuristic, $\text{range}(\mathbf{Y})/\max(\mathbf{Y})$, clipped to $[0.05,0.15]$ for stability.

\textbf{Acquisition and Batch Selection (sequential greedy).}
UCB acquisition is $\alpha(\mathbf{x}) = \mu(\mathbf{x}) + \beta \cdot \sigma(\mathbf{x})$, with $\beta$ decaying from 3.5 to 1.3 over 15 iterations.
Batch selection uses sequential greedy selection: points are picked one-by-one, and after each pick we update a diversity penalty (repulsion) before selecting the next point.
This is a fast approximation to joint batch optimization and is appropriate under the runtime constraint.
After each selection, remaining candidates are penalized by $\exp(-d_{\min}/r)$ where $r$ shrinks from 0.22 to 0.08; penalty strength is 0.45.
Cell types not yet in the current batch additionally receive a 5\% UCB bonus.

\textbf{Stagnation Detection.}
If no improvement occurs for 3+ iterations, $\beta$ is boosted by 1.4$\times$ (capped at 4.0).
If one cell type comprises $>$85\% of recent samples, $\beta$ is boosted by 1.2$\times$.
These mechanisms prevent premature convergence and encourage exploration when needed.

\section{Pseudocode}

\begin{algorithm}[H]
\caption{Robust Batch Bayesian Optimization}
\begin{algorithmic}[1]
\Require Bounds $\mathcal{B}$, cell types $\mathcal{C}$, iterations $T=15$, batch size $m=5$
\Ensure Best parameters $\mathbf{x}^*$, best titre $f^*$
\State $\mathbf{X} \gets \textsc{SobolDesign}(6, \mathcal{B}, \mathcal{C})$; \quad $\mathbf{Y} \gets \textsc{Evaluate}(\mathbf{X})$
\State $\text{no\_improve} \gets 0$
\For{$t = 1, \ldots, T$}
    \If{elapsed $> 52$s} \textbf{break} \EndIf
    \State Estimate noise using normalized range:
    \State \hspace{1em} $\sigma_n \gets \text{clip}(0.08 \cdot \text{range}(\mathbf{Y})/\max(\mathbf{Y}), 0.05, 0.15)$
    \State $\text{GP} \gets \textsc{MixedGP}(\mathbf{X}, \mathbf{Y}, \theta_{\text{cat}}=0.25, \sigma_n)$ \Comment{Adaptive length scales}
    \State $\beta \gets \max(1.3, \; 3.5 - 2.2 \cdot (t-1)/14)$
    \State $r \gets \max(0.08, \; 0.22 - 0.14 \cdot (t-1)/14)$
    \If{$\text{no\_improve} \geq 3$} $\beta \gets \min(4.0, 1.4 \cdot \beta)$ \EndIf
    \If{cell type dominance $> 85\%$ in last 10 samples} $\beta \gets \min(4.0, 1.2 \cdot \beta)$ \EndIf
    \State $n_g \gets \{12000, 8000, 5000\}$ for $t \in \{1\text{--}3, 4\text{--}6, 7\text{--}15\}$
    \State $\mathbf{C} \gets \textsc{SobolCandidates}(n_g, \mathcal{B}, \mathcal{C})$
    \If{$t \geq 3$}
        \State $n_\ell \gets 2500$ if $t < 6$ else $4000$
        \State $\mathbf{C} \gets \mathbf{C} \cup \textsc{LocalCandidates}(\arg\max \mathbf{Y}, n_\ell, \text{window}=12\%, \text{min}=3\%)$
        \Comment{70\% best type, 30\% others}
    \EndIf
    \State $\mathbf{B} \gets []$; $\mathbf{B}_{\text{cats}} \gets []$
    \For{$j = 1, \ldots, m$} \Comment{Sequential greedy batch construction}
        \State $\alpha_i \gets \mu(\mathbf{c}_i) + \beta \cdot \sigma(\mathbf{c}_i)$ for all $\mathbf{c}_i \in \mathbf{C}$
        \If{$\mathbf{B} \neq \emptyset$}
            \State $d_i \gets \min_{\mathbf{b} \in \mathbf{B}} \|\mathbf{c}_i - \mathbf{b}\|$; \quad $\alpha_i \gets \alpha_i \cdot (1 - 0.45 \cdot e^{-d_i/r})$
            \State $\alpha_i \gets \alpha_i + 0.05 \cdot |\bar{\alpha}|$ if cell type not in $\mathbf{B}_{\text{cats}}$
        \EndIf
        \State $\mathbf{B} \gets \mathbf{B} \cup \{\arg\max_i \alpha_i\}$; update $\mathbf{B}_{\text{cats}}$; remove from $\mathbf{C}$
    \EndFor
    \State $\mathbf{Y}_B \gets \textsc{Evaluate}(\mathbf{B})$; \quad $\mathbf{X} \gets \mathbf{X} \cup \mathbf{B}$; \quad $\mathbf{Y} \gets \mathbf{Y} \cup \mathbf{Y}_B$
    \State $\text{no\_improve} \gets 0$ \textbf{if} $\max \mathbf{Y}$ improved \textbf{else} $\text{no\_improve} + 1$
\EndFor
\State \Return $\arg\max \mathbf{Y}$, $\max \mathbf{Y}$
\end{algorithmic}
\end{algorithm}

\vspace{0.3em}
\noindent\textbf{Summary.}
Sobol sequences ensure uniform initial coverage.
The mixed GP with adaptive length scales and $\theta_{\text{cat}}=0.25$ balances knowledge transfer with cell-type distinction.
Decaying $\beta$ (3.5$\to$1.3) shifts from exploration to exploitation.
Cross-cell-type local search uses a 12\% window (min 3\%) to refine promising regions without collapsing the search area.
Sequential greedy batch construction plus repulsion and cell-type bonuses maintain diversity.
Stagnation detection with $\beta$-boosting helps escape local optima.

\end{document}
